\clearpage
\begin{proofbox}
\textbf{\textcolor{CProof}{方法一：双曲线定义＋平方差（考场速证）}}

设 $d_1=|PF_1|,d_2=|PF_2|$，由定义有 $|d_1-d_2|=2a$。
两点距离平方相减，共同的 $y_0^2$ 消去：
\[
d_1^2-d_2^2=[(x_0+c)^2+y_0^2]-[(x_0-c)^2+y_0^2]=4cx_0.
\]
因为 $d_1+d_2>0$，平方差的符号就是 $d_1-d_2$ 的符号，故能严格判断哪段较长。

\textbf{右支 $x_0\ge a$：} 平方差为正，所以 $d_1>d_2$，$d_1-d_2=2a$。
\[
d_1+d_2=\frac{4cx_0}{2a}=2ex_0.
\]
由和差求两数：
\[
\boxed{d_1=ex_0+a},\qquad\boxed{d_2=ex_0-a}.
\]
\textbf{左支 $x_0\le-a$：} 平方差为负，所以 $d_1<d_2$，$d_1-d_2=-2a$。
\[
d_1+d_2=\frac{4cx_0}{-2a}=-2ex_0.
\]
于是
\[
\boxed{d_1=-ex_0-a},\qquad\boxed{d_2=-ex_0+a}.
\]
两支统一为 $\boxed{d_1=|a+ex_0|,\ d_2=|a-ex_0|}$。
本法只用定义、两点距离平方差和加减消元，不必代入双曲线方程展开。
\end{proofbox}
\clearpage
\begin{proofbox}
\textbf{\textcolor{CProof}{方法二：先设两长度的平均值 $t$（首推理解）}}

两段差为 $2a$，所以让长、短分别为 $t+a,t-a$，只用一个未知量。
\begin{center}
\begin{tikzpicture}[x=1cm,y=1cm,line cap=round,line join=round,font=\small]
\draw[black!25,rounded corners=3pt] (0,0) rectangle (14.4,8.6);
\node[font=\small\bfseries] at (7.2,8.1) {把长短两段对齐：平均长度 $t$ 的两侧各差一个 $a$};
\node at (7.2,7.35) {平移到同一起点，只比较长度；示例 $a=2,\ t=9/2$。};
\fill[blue!40] (2,6.0) rectangle (11.75,6.2);
\node[blue!70!black,above] at (5.3,6.2) {较长 $t+a=13/2$};
\draw[black!65,dashed,line width=1.2pt] (2,4.6)--(8.75,4.6);
\node[above] at (5.3,4.6) {平均长度 $t=9/2$};
\fill[orange!50] (2,2.9) rectangle (5.75,3.1);
\node[orange!85!black,above] at (3.9,3.1) {较短 $t-a=5/2$};
\draw[black!30,dotted] (2,2.5)--(2,6.6);
\draw[purple!65,<->] (8.75,5.35)--(11.75,5.35) node[midway,above] {$a=2$};
\draw[purple!65,<->] (5.75,3.8)--(8.75,3.8) node[midway,above] {$a=2$};
\draw[black!60,<->] (5.75,1.8)--(11.75,1.8) node[midway,below] {$2a=4$};
\draw[black!30,dotted] (5.75,1.7)--(5.75,2.9) (11.75,1.7)--(11.75,6.0);
\node at (7.2,.75) {$t=\dfrac{\text{长}+\text{短}}2=e|x_0|$；$t$ 是两个长度的平均值，不是焦点连线的中点。};
\end{tikzpicture}
\end{center}

\textbf{右支：} $d_1=t+a,d_2=t-a$，平方相减有
\[
(t+a)^2-(t-a)^2=4at=4cx_0\ \Longrightarrow\ t=ex_0.
\]
\textbf{左支：} $d_1=t-a,d_2=t+a$，同理
\[
(t-a)^2-(t+a)^2=-4at=4cx_0\ \Longrightarrow\ t=-ex_0.
\]
因此两支都有 $t=e|x_0|$，代回即得“同侧短、异侧长”。
图中的三个长度从同一起点对齐，并非在原焦点三角形中取一个几何中点。
\end{proofbox}
\clearpage
\begin{proofbox}
\textbf{\textcolor{CProof}{方法三：标准方程消元、配方、开方}}

点在双曲线上，故 $y_0^2=\frac{b^2}{a^2}(x_0^2-a^2)$，且 $c^2=a^2+b^2$。
\[
\begin{aligned}
d_1^2&=(x_0+c)^2+y_0^2\\
&=x_0^2+2cx_0+c^2+\frac{b^2}{a^2}x_0^2-b^2\\
&=\frac{c^2}{a^2}x_0^2+2cx_0+a^2=(ex_0+a)^2.
\end{aligned}
\]
把 $c$ 换成 $-c$，同理得 $d_2^2=(ex_0-a)^2$。开方必须保留绝对值：
\[
\boxed{d_1=|ex_0+a|=|a+ex_0|},\qquad
\boxed{d_2=|ex_0-a|=|a-ex_0|}.
\]
\textbf{为什么不能像椭圆那样直接去掉绝对值？}
因为 $e>1,|x_0|\ge a$；右支 $ex_0-a\ge c-a>0$，
左支 $ex_0+a\le a-c<0$。去绝对值前必须先判断所在支。
易错页的数轴图把两个括号的符号一次列清。

\textbf{顶点也成立：} $P=(a,0)$ 时 $(d_1,d_2)=(c+a,c-a)$；
$P=(-a,0)$ 时两者对调。焦半径均非零；仅焦点三角形退化。
\end{proofbox}
\clearpage
\begin{proofbox}
\textbf{\textcolor{CProof}{方法四：焦点—准线第二定义（选学）}}

\textbf{前提：} 已学过并证明第二定义 $|PF|/d(P,\ell)=e$；此处只应用它，不用它循环证明自身。
图中 $A_-,A_+$ 分别是左右顶点，准线比顶点更靠近原点。
\begin{center}
\begin{tikzpicture}[x=1cm,y=1cm,line cap=round,line join=round,font=\small]
\draw[black!25,rounded corners=3pt] (0,0) rectangle (14.4,10);
\node[font=\small\bfseries] at (7.2,9.5) {双曲线准线在两顶点之间；右支点位于两准线右侧};
\begin{scope}[shift={(5.1,4.5)},x=0.9cm,y=0.9cm]
\draw[black!30,->] (-4.4,0)--(4.5,0) node[right] {$x$};
\draw[black!65,thick,domain=-3.5:3.5,samples=100,variable=\z] plot ({-1*2*sqrt(1+\z*\z/5)},{\z});
\draw[black!65,thick,domain=-3.5:3.5,samples=100,variable=\z] plot ({1*2*sqrt(1+\z*\z/5)},{\z});
\coordinate (P) at (3,2.5);
\coordinate (Fone) at (-3,0);
\coordinate (Ftwo) at (3,0);
\draw[blue!70!black,line width=1.1pt] (Fone)--(P);
\draw[orange!85!black,line width=1.1pt] (Ftwo)--(P);
\foreach \q in {P,Fone,Ftwo} \fill (\q) circle (1.6pt);
\draw[purple!65,dashed] (-1.333333,-3.3)--(-1.333333,3.7);
\draw[purple!65,dashed] (1.333333,-3.3)--(1.333333,3.7);
\coordinate (Kone) at (-1.333333,2.5); \coordinate (Ktwo) at (1.333333,2.5);
\draw[blue!70!black,dashed,line width=2pt] (Kone)--(P);
\draw[orange!85!black,densely dashed,line width=.8pt] (Ktwo)--(P);
\draw[purple!65] (-1.333333,2.25)--(-1.083333,2.25)--(-1.083333,2.5);
\draw[purple!65] (1.333333,2.25)--(1.583333,2.25)--(1.583333,2.5);
\node[below left] at (Kone) {$K_1$}; \node[below left] at (Ktwo) {$K_2$};
\node[below right] at (P) {$P$};
\node[below] at (Fone) {$F_1$}; \node[below] at (Ftwo) {$F_2$};
\fill (-2,0) circle (1.5pt); \draw[black!40] (-2,0)--(-2.6,-.75) node[below] {$A_-$};
\fill (2,0) circle (1.5pt); \draw[black!40] (2,0)--(2.6,-.75) node[below] {$A_+$};
\end{scope}
\node[purple!65] at (3.9,8.55) {$x=-\dfrac ae$};
\node[purple!65] at (6.3,8.55) {$x=\dfrac ae$};
\node[blue!70!black] at (4.95,7.2) {$x_0+\dfrac ae$};
\node[orange!85!black] at (7.2,7.2) {$x_0-\dfrac ae$};
\node[align=left,anchor=north west] at (10,7.7) {$d_1=e\,PK_1$\\[1em]$d_2=e\,PK_2$\\[1.3em]右支：\\[.6em]$d_1=ex_0+a$\\[.6em]$d_2=ex_0-a$};
\node at (7.2,.55) {两条水平距离有重合部分，用蓝粗虚线、橙细虚线区分；左支取距离绝对值。};
\end{tikzpicture}
\end{center}

与 $F_1,F_2$ 对应的准线依次是 $x=-a/e,x=a/e$，所以
\[
\begin{aligned}
d_1&=e\left|x_0+\frac ae\right|=|ex_0+a|,\\
d_2&=e\left|x_0-\frac ae\right|=|ex_0-a|.
\end{aligned}
\]
右支点位于两条准线右侧，左支点位于两条准线左侧，故水平距离的去绝对值形式不同。
没有学过第二定义时，使用方法一或方法二即可。
\end{proofbox}
